\input{../tex-inputs/latex-leseansicht-vorspann}

\section[Gustav Schwarzkopf an Arthur Schnitzler, 7. 12. 1899]{L04288 Gustav Schwarzkopf an Arthur Schnitzler, 7. 12. 1899}
\nopagebreak\mylabel{L04288v}
\rehead{ }\normalsize\beginnumbering\briefempfaengerindex{Schnitzler, Arthur@\textsc{Schnitzler, Arthur}!zzzSchwarzkopf, Gustav@\emph{von Gustav Schwarzkopf}!1899-12-072@{7. 12. 1899}|(be}
\toendnotes[C]{\smallbreak\pagebreak[2]}
\correspDesc{Versand  durch Gustav Schwarzkopf am 7. 12. 1899 in Wien
\newline{}Übermittlung  am 7. 12. 1899 in Wien
\newline{}Zustellung  am 8. 12. 1899 in Wien
\newline{}Erhalt  durch Arthur Schnitzler am 8. 12. 1899 in Wien}\toendnotes[C]{\smallbreak}
\Standort{CUL, Schnitzler, B 96a.}
\physDesc{Postkarte, 223 Zeichen
\newline{}Handschrift: schwarze Tinte, deutsche Kurrent
\newline{}Versand: 1) Stempel: »\nobreak{}\oindex{I., Innere Stadt@\textbf{I., Innere Stadt}, \emph{Verwaltungsgebiet}|pwk}1/1 Wien 1, 7 12 99, 10–11V\nobreak{}«.   2) Stempel: »\nobreak{}\oindex{IX., Alsergrund@\textbf{IX., Alsergrund}, \emph{Verwaltungsgebiet}|pwk}Wien 72, 8. 12. 99, 8.V, \textcolor{gray}{Bestellt}\nobreak{}«. }\pstart{}{\pb}Herrn \textsc{D\textsuperscript{r} Arthur Schnitzler}\pend{}\pstart{}\textsc{Wien\oindex{Wien@\textbf{Wien}, \emph{Verwaltungsgebiet}|pw}}\pend{}\pstart{}IX Frankgaſſe 1\oindex{Wien@\textbf{Wien}!IX., Alsergrund@\textbf{IX., Alsergrund}!Frankgasse 1@\textbf{Frankgasse 1}, \emph{Wohngebäude}|pw}.\pend{}{\bigskip}\vspace{1em}
\pstart
           \raggedleft{}{\pb}7. XII 99\pend
           \vspace{0.5em}
\pstart
           Lieber Arthur, ich bin Freitag bei \textsc{Chiavacci\pwindex{Chiavacci, Vincenz 15.\,6.\,1847 Wien – 2.\,2.\,1916 ebd.@\textsc{Chiavacci, Vincenz} (15.\,6.\,1847 Wien – 2.\,2.\,1916 ebd.), \emph{Schriftsteller, Journalist, Beamter}|pw}} geladen, großes Zauberfeſt zu EhrenGanghofers\pwindex{Ganghofer, Ludwig 7.\,7.\,1855 Kaufbeuren – 24.\,7.\,1920 Tegernsee@\textsc{Ganghofer, Ludwig} (7.\,7.\,1855 Kaufbeuren – 24.\,7.\,1920 Tegernsee), \emph{Schriftsteller}|pw}\pwindex{Ganghofer, Katharina 7.\,7.\,1856 Pest – 8.\,4.\,1930 München@\textsc{Ganghofer, Katharina} (7.\,7.\,1856 Pest – 8.\,4.\,1930 München), \emph{Schauspielerin}|pw}), ka{\geminationn} alſo von Ihrer liebenswürdigen
               Einladung – u. ſ\textcolor{gray}{o} w.\pend
           
\pstart
           Ich grüße Sie herzlich{\\[\baselineskip]}\spacefill\mbox{Gustav}\pend
           \leftskip=0em{}\selectlanguage{ngerman}\endnumbering\briefempfaengerindex{Schnitzler, Arthur@\textsc{Schnitzler, Arthur}!zzzSchwarzkopf, Gustav@\emph{von Gustav Schwarzkopf}!1899-12-072@{7. 12. 1899}|)be}\mylabel{L04288h}
\begin{anhang}
\end{anhang}\newcommand{\dateiname}{L04288}\newcommand{\titel}{Gustav Schwarzkopf an Arthur Schnitzler, 7. 12. 1899}\newcommand{\editorInnen}{Herausgegeben von Jahnke, SelmaMüller, Martin Anton}\input{../tex-inputs/latex-leseansicht-abspann}
